\documentclass{article}
\usepackage{amsmath, amssymb, amsthm}
\usepackage{booktabs}
\usepackage{geometry}
\usepackage{hyperref}
\usepackage{cleveref}

\geometry{margin=1in}
\hypersetup{colorlinks=true, urlcolor=blue}

% Shared macros for the multi-mode mixed-criticality scheduling papers.
% % Shared macros for the multi-mode mixed-criticality scheduling papers.
% % Shared macros for the multi-mode mixed-criticality scheduling papers.
% \input{macros} from any document in this directory rather than redefining
% these locally -- a macro defined in one paper's preamble and silently
% missing from a sibling's is exactly the class of error that put \round and
% \E undefined in evaluation_methodology.tex.

\newcommand{\N}{\mathbb{N}}
\newcommand{\R}{\mathbb{R}}
\newcommand{\Z}{\mathbb{Z}}
\newcommand{\E}{\mathbb{E}}
\newcommand{\Var}{\mathrm{Var}}
\newcommand{\Prob}{\mathbb{P}}
\newcommand{\Ind}{\mathbf{1}}
\newcommand{\ceil}[1]{\lceil #1 \rceil}
\newcommand{\floor}[1]{\lfloor #1 \rfloor}
\newcommand{\round}[1]{\lfloor #1 \rceil}
 from any document in this directory rather than redefining
% these locally -- a macro defined in one paper's preamble and silently
% missing from a sibling's is exactly the class of error that put \round and
% \E undefined in evaluation_methodology.tex.

\newcommand{\N}{\mathbb{N}}
\newcommand{\R}{\mathbb{R}}
\newcommand{\Z}{\mathbb{Z}}
\newcommand{\E}{\mathbb{E}}
\newcommand{\Var}{\mathrm{Var}}
\newcommand{\Prob}{\mathbb{P}}
\newcommand{\Ind}{\mathbf{1}}
\newcommand{\ceil}[1]{\lceil #1 \rceil}
\newcommand{\floor}[1]{\lfloor #1 \rfloor}
\newcommand{\round}[1]{\lfloor #1 \rceil}
 from any document in this directory rather than redefining
% these locally -- a macro defined in one paper's preamble and silently
% missing from a sibling's is exactly the class of error that put \round and
% \E undefined in evaluation_methodology.tex.

\newcommand{\N}{\mathbb{N}}
\newcommand{\R}{\mathbb{R}}
\newcommand{\Z}{\mathbb{Z}}
\newcommand{\E}{\mathbb{E}}
\newcommand{\Var}{\mathrm{Var}}
\newcommand{\Prob}{\mathbb{P}}
\newcommand{\Ind}{\mathbf{1}}
\newcommand{\ceil}[1]{\lceil #1 \rceil}
\newcommand{\floor}[1]{\lfloor #1 \rfloor}
\newcommand{\round}[1]{\lfloor #1 \rceil}


\title{Choosing an Optimisation Method for Severity-Ladder Design}
\author{}
\date{}

\begin{document}

\maketitle

\section{Introduction}

Phase 3 of this work asks: for $k$ degradation levels, what severities
$\chi_2, \dots, \chi_{k-1}$ should trigger levels $L_2, \dots, L_{k-1}$
(severity-indexed response time and admissibility, task\_model.tex), and what
optimisation method should be used to find them? The method cannot be chosen from habit or
convenience -- the wrong choice either wastes evaluation budget the problem
does not need spent, or under-samples a domain that needs full coverage. This
document characterises the search problem first (\Cref{sec:landscape}), then
uses that characterisation to assess the candidate optimisation techniques
(\Cref{sec:related-work}) and justify the one selected
(\Cref{sec:recommendation}).

\section{Problem Characterisation: The Optimisation Landscape}
\label{sec:landscape}

\subsection{Decision Variables and Domain}

A $k$-level severity ladder has $k - 1$ trigger points, but the deepest,
$\chi_{k-1}$, is not free: the safety argument (admissible drop sets,
task\_model.tex) requires $\chi_1 = 0$, and pinning $\chi_{k-1}$ recovers the
classic two-level guarantee exactly at the deepest level. The optimisation is therefore over $\mathbf{k - 2}$ free variables
$\chi_2 \leq \chi_3 \leq \dots \leq \chi_{k-2}$, each bounded in $(0, 1]$ and
subject to a total order. This is small in every case this work considers:

\begin{table}[h]
\centering
\caption{Free dimensions and feasible-set size by level count}
\begin{tabular}{cccc}
\toprule
$k$ & free dimensions & unordered hypercube ($g=20$) & ordered feasible set \\
\midrule
3 & 1 & 400 & \textbf{20} \\
4 & 2 & 8{,}000 & \textbf{210} \\
5 & 3 & 160{,}000 & \textbf{1{,}540} \\
\bottomrule
\end{tabular}
\end{table}

\noindent The right-hand column is not the hypercube $g^{k-2}$ but the number
of non-decreasing $(k{-}2)$-tuples drawn from a grid of $g$ values,
$\binom{g+k-3}{k-2}$ (a standard multiset/"stars-and-bars" count): the ordering
constraint is not an incidental detail but the dominant factor in how large
the space actually is. A search method that samples the unconstrained box and
repairs infeasible points by sorting -- the natural approach for a
population-based or surrogate-based method operating on a free vector --
wastes evaluations on points that collapse onto the same feasible tuple, and
in the worst case biases the sampling measure over the feasible region without
the practitioner intending it to.

\subsection{The Objective is a Noisy, Expensive-per-Sample Simulator Output}

$\Phi$ (metrics\_objective.md) is not observed directly; it is estimated from
a discrete-event simulation whose execution times and HI-criticality
occurrences are drawn at random. Per-seed coefficients of variation, measured
on this simulator: $0.157$ (NiD), $0.314$ (JNE), $0.199$ (degraded ticks).
Neither the objective nor its gradient is available in closed form, so the
landscape is a black box in the sense every method below has to contend with,
not merely difficult to solve analytically.

A single simulation is cheap ($\approx 43\,\mathrm{ms}$ for $10^6$ ticks) and
embarrassingly parallel across seeds and task sets, but the noise above means
one evaluation at one lattice point is not one sample -- it is a paired mean
over a seed block, and the block has to be large enough to resolve the effect
size the study cares about (\S\ref{sec:related-work} returns to what this
implies for methods that spend their budget on \emph{many} points rather than
\emph{repeated} samples at few).

\subsection{What Is Not Yet Known}

Whether $\Phi$ is unimodal, smooth, or has a plateau over the severity domain
is an open empirical question, not an assumption available to lean on. A
method whose efficiency depends on such structure -- surrogate modelling,
gradient estimation, population convergence -- is therefore choosing to bet on
a property this document cannot yet certify. A method that does not depend on
it is not making that bet, at the cost of not exploiting the structure if it
does turn out to hold. This tension is the crux of \Cref{sec:related-work}.

\section{Related Work: Optimisation Techniques for Expensive Stochastic Objectives}
\label{sec:related-work}

Optimising the parameters of a stochastic simulation, rather than a
closed-form function, is the subject of \emph{simulation optimisation}, whose
central tension Fu~\cite{fu2002optimization} frames precisely: theoretical
convergence results are proven for algorithms that are rarely what
practitioners actually run, while the general-purpose metaheuristics that are
widely used borrow convergence properties from deterministic optimisation that
do not transfer cleanly to a noisy objective. Four families of technique are
relevant to the problem in \Cref{sec:landscape}.

\subsection{Exhaustive Search and Design of Experiments}

Design of Experiments (DOE) fits a statistical model -- typically linear or
low-order polynomial -- to a systematically chosen sample of the parameter
space, trading exhaustive coverage for a smaller, structured design (e.g.\ a
factorial or Latin-hypercube design) when the space is too large to enumerate.
Tate et al.~\cite{R:Tate:2009a} apply exactly this trade-off to a
directly comparable problem: tuning multiple sensornet protocol parameters
against several competing performance metrics, a genuinely high-dimensional,
combinatorially explosive space where full enumeration is not an option. Their
DOE fit yielded a more generally applicable model of the parameter space; a
competing evolutionary search (\Cref{sec:evolutionary}) found a broader range
of viable solutions at lower cost for the specific instance evaluated. That
finding does not transfer here, and the reason it does not is itself
informative: their conclusion favours evolutionary search precisely where full
enumeration is infeasible, and \Cref{sec:landscape} establishes that the
severity ladder is not in that regime -- 1{,}540 points at $k=5$ is exhaustively
enumerable, so DOE's usual purpose, choosing which subset of a space to sample
because the whole space cannot be, does not arise. What remains applicable is
the weaker, unconditional claim: over a space small enough to enumerate,
exhaustive search dominates any sampling design of it, since a sample is only
ever a strict subset of what enumeration already covers, at correspondingly
less information per unit of the region actually examined.

\subsection{Evolutionary and Population-Based Metaheuristics}
\label{sec:evolutionary}

Genetic algorithms~\cite{Gordon93serialand, WHITELY93} and related
population-based metaheuristics such as simulated
annealing~\cite{kirkpatrick83optimization} search by iteratively
perturbing and selecting among a working set of candidate solutions, and are
the default choice when a space is too large or too irregular to enumerate and
no exploitable structure is known in advance -- exactly the case in
Tate et al.'s protocol-tuning problem above. Applied here, the same
population-based iteration that gives these methods their reach becomes a
liability: each generation adds evaluations, and \Cref{sec:landscape}'s
per-seed noise means every added evaluation is a further noisy draw. Selecting
the best of $M$ such draws is optimistically biased by
$\sigma \sqrt{2 \ln M}$ (a standard result for the expected maximum of
$M$ noisy samples). A population of 50 run for 100 generations -- an
unremarkable configuration -- makes 5{,}000 evaluations; at this simulator's
measured noise, that inflates the reported optimum's bias to roughly
$12\%$, more than double the $5\%$ practical-significance threshold this work
uses to call a result meaningful (metrics\_objective.md). \emph{A genetic
algorithm run against a perfectly flat objective would return a
publishable-looking optimum.} This is a general property of any method whose
evaluation count scales with iterations rather than with the size of the space
being covered, and it argues against every method in this family for a
domain this small, independent of how well-suited any of them are in a
domain too large to enumerate.

\subsection{Bayesian Optimisation and Surrogate Models}

Bayesian optimisation~\cite{jones1998efficient, shahriari2016taking} fits a
probabilistic surrogate -- typically a Gaussian process -- to the evaluations
seen so far, and chooses the next point by an acquisition function that
balances exploiting the surrogate's optimum against exploring where it is most
uncertain. It is the standard choice when each evaluation is expensive enough
that minimising the \emph{number} of evaluations dominates every other
concern, and its sample efficiency is real precisely under that condition. It
is not the condition here: \Cref{sec:landscape} costs a single simulation at
tens of milliseconds and the whole lattice at a few CPU-hours, so there is no
evaluation budget for a surrogate to conserve. Two further mismatches follow
from \Cref{sec:landscape} specifically. First, the deliverable is a response
surface, not an optimum -- Bayesian optimisation is built to concentrate
samples near a believed optimum, which is the opposite of what characterising
a surface requires, and recovering the surface afterwards needs a further,
separately budgeted sampling pass. Second, the surrogate assumes a smoothness
the objective has not yet been shown to have (\Cref{sec:landscape}); a
Gaussian-process surrogate additionally needs an explicit noise term fitted
correctly to the measured heteroscedastic variance above, which is one more
unverified assumption riding on an already-unverified one.

\subsection{Gradient-Based and Derivative-Free Local Search}

Classical gradient-based optimisation is inapplicable outright: $\Phi$ is the
output of a discrete-event simulation with no closed form and no analytic
derivative, and finite-difference gradient estimates would themselves be noisy
at the measured coefficients of variation, compounding rather than resolving
the problem \Cref{sec:evolutionary} raises. Derivative-free local search
(pattern search, Nelder--Mead) removes the need for a derivative but keeps the
core liability: it is iterative and local, so it inherits both the winner's
curse of repeated noisy evaluation and a dependence on the smoothness this
document has not established.

\section{Recommendation}
\label{sec:recommendation}

\textbf{Exhaustive enumeration of the ordered severity lattice
(\Cref{sec:landscape}), evaluated under common random numbers.} Every
technique surveyed in \Cref{sec:related-work} earns its place by trading
coverage for sample efficiency in a large or unstructured space; the severity
ladder's space is neither. DOE's usual justification -- covering a space too
large to enumerate -- does not apply once the ordering constraint is
respected directly rather than repaired after the fact. Evolutionary and
Bayesian methods both spend evaluations to save evaluations, at a specific,
quantified cost (winner's-curse bias, or an unrepresentative sample of the
surface) that this domain does not need paid. Enumeration additionally
reports the whole surface natively, has one hyperparameter -- grid resolution
-- against a genetic algorithm's at least eight or a Gaussian process's
comparable number, and needs no additional dependency. This recommendation is
conditional, not permanent: it should be revisited if $k$ grows past roughly
six levels, or if the trigger vector gains dimensions (per-task rather than
per-level triggers scale with $n$, not $k$, and enumeration stops being
affordable) -- at that point a Bayesian method adopted \emph{with} the paired
evaluation protocol below is the natural replacement, since the case against
it here is about cost, not principle.

Whichever method is used, it must respect the noise floor
\Cref{sec:landscape} establishes: configurations are compared as paired
samples under common random numbers, never as independent means, since an
unpaired comparison at this simulator's measured noise needs roughly two
orders of magnitude more task sets to resolve the same effect
(metrics\_objective.md, evaluation\_methodology.tex \S3). This is the one
requirement every method in \Cref{sec:related-work} shares, and the one a
choice of \emph{method} cannot substitute for.

\bibliographystyle{plain}
\bibliography{project_bib,references}

\end{document}
